\documentclass[11pt,a4paper]{article}
\usepackage[utf8]{inputenc}
\usepackage[vietnamese]{babel}
\usepackage{amsmath,amssymb}
\usepackage{geometry}
\geometry{a4paper, total={180mm,260mm}, left=15mm, right=15mm, top=18mm, bottom=18mm}

\usepackage{tikz}
\usetikzlibrary{shapes.geometric, arrows.meta, positioning, calc}
\tikzset{
  arr/.style={-{Stealth[scale=0.8]}},
  every node/.append style={align=center}
}
\usepackage{tcolorbox}
\tcbuselibrary{skins,breakable}
\usepackage{fontawesome}
\usepackage{fancyhdr}
\usepackage{enumitem}

\pagestyle{fancy}
\fancyhf{}
\renewcommand{\headrulewidth}{0.5pt}
\renewcommand{\footrulewidth}{0.5pt}
\fancyhead[L]{\small\bfseries\color{violet!85!black}\faBook\ SỔ TAY KIẾN THỨC TOÁN 2}
\fancyhead[R]{\small\bfseries\color{violet!85!black}Lớp Toán Thầy Lĩnh -- loptoanthaylinh.vn}
\fancyfoot[C]{\small\color{gray}-- Sổ tay kiến thức: Bài 14. Luyện tập chung Chủ đề 2 --}

\setlength{\parindent}{0pt}
\setlength{\parskip}{0pt}

\begin{document}

% --- BANNER BÀI HỌC ---
\begin{tcolorbox}[
  enhanced,
  colback=violet!8!white,
  colframe=violet!70!black,
  arc=3mm, boxrule=1.2pt,
  left=3mm, right=3mm, top=2mm, bottom=2mm,
  halign=center
]
  {\LARGE\bfseries\color{violet!90!black}\faBookmark\ BÀI 14: LUYỆN TẬP CHUNG CHỦ ĐỀ 2}\\[1mm]
  {\itshape\footnotesize\color{black!75} Đại tổng kết phép cộng, phép trừ có nhớ trong phạm vi 20 và các bài toán thực tế trên 1 trang A4}
\end{tcolorbox}

\vspace{-2mm}

\noindent
\begin{minipage}[t]{0.49\textwidth}\vspace{0pt}
\begin{tcolorbox}[
  enhanced,
  colback=violet!4!white,
  colframe=violet!75!black,
  arc=2mm, boxrule=0.8pt,
  left=2mm, right=2mm, top=1.8mm, bottom=1.8mm,
  title={\footnotesize\bfseries\color{white}\faKey\ 1. Quy tắc tính nhẩm cốt lõi}
]
{\footnotesize
\begin{itemize}\setlength{\itemsep}{1pt}\setlength{\leftskip}{-2.5mm}
  \item \textbf{Phép cộng qua 10:}\\
  Tách số hạng sau để làm tròn $10$, rồi cộng phần còn lại với $10$.\\
  $9+c$ (tách $1$), $8+c$ (tách $2$), $7+c$ (tách $3$), $6+c$ (tách $4$).
  \item \textbf{Phép trừ qua 10:}\\
  Trừ để được $10$ rồi trừ tiếp số còn lại, hoặc lấy $10$ trừ số trừ rồi cộng đơn vị.
  \item \textbf{Dãy tính 3 số:}\\
  Tính tuần tự từ trái sang phải: $9 + 5 - 6 = 14 - 6 = 8$.
\end{itemize}
}
\end{tcolorbox}
\end{minipage}
\hfill
\begin{minipage}[t]{0.49\textwidth}\vspace{0pt}
\begin{tcolorbox}[
  enhanced,
  colback=teal!4!white,
  colframe=teal!75!black,
  arc=2mm, boxrule=0.8pt,
  left=2mm, right=2mm, top=1.8mm, bottom=1.8mm,
  title={\footnotesize\bfseries\color{white}\faSliders\ 2. Mối liên hệ Phép cộng - Phép trừ}
]
{\footnotesize
\begin{itemize}\setlength{\itemsep}{1pt}\setlength{\leftskip}{-2.5mm}
  \item Từ một phép tính cộng, ta luôn viết được \textbf{hai phép tính trừ tương ứng}:
  $$a + b = c \implies c - a = b \text{ và } c - b = a$$
  \textit{Ví dụ:} $9 + 7 = 16 \implies 16 - 9 = 7$ và $16 - 7 = 9$.
  \item \textbf{Tìm thành phần chưa biết:}\\
  $\square + 8 = 15 \implies \square = 15 - 8 = 7$.\\
  $14 - \square = 8 \implies \square = 14 - 8 = 6$.
\end{itemize}
}
\end{tcolorbox}
\end{minipage}

\vspace{-1.5mm}
\begin{tcolorbox}[
  enhanced,
  colback=orange!4!white,
  colframe=orange!80!black,
  arc=2mm, boxrule=0.8pt,
  left=2mm, right=2mm, top=1.5mm, bottom=1.5mm,
  title={\footnotesize\bfseries\color{white}\faSitemap\ 3. Sơ đồ tư duy 4 nhánh: Tổng kết Chủ đề 2 (Cộng, trừ qua 10)}
]
\centering
\begin{tikzpicture}[
  node distance=0.8cm,
  box/.style={rectangle, rounded corners=2pt, draw=violet!80!black, fill=violet!15, text width=2.4cm, align=center, font=\scriptsize\bfseries, inner sep=3pt},
  branch/.style={rectangle, rounded corners=2pt, draw=#1!80!black, fill=#1!10, text width=2.5cm, align=center, font=\scriptsize, inner sep=2.5pt},
  arr/.style={-{Stealth[scale=0.8]}, thick, #1!80!black}
]
  \node[box] (root) {CHỦ ĐỀ 2\\CỘNG, TRỪ\\PHẠM VI 20};
  \node[branch=teal, above left=0.3cm and 0.4cm of root] (b1) {\textbf{CỘNG QUA 10}\\Tách gộp đủ 10\\Cộng phần còn lại};
  \node[branch=orange, above right=0.3cm and 0.4cm of root] (b2) {\textbf{TRỪ QUA 10}\\Trừ về 10 rồi trừ tiếp\\$10 - b$ cộng đơn vị};
  \node[branch=green, below left=0.3cm and 0.4cm of root] (b3) {\textbf{THÊM -- BỚT}\\Thêm vào $\implies (+)$\\Bớt đi $\implies (-)$};
  \node[branch=purple, below right=0.3cm and 0.4cm of root] (b4) {\textbf{NHIỀU HƠN -- ÍT HƠN}\\Tìm số lớn $\implies (+)$\\Tìm số bé $\implies (-)$};

  \draw[arr=teal] (root) -- (b1);
  \draw[arr=orange] (root) -- (b2);
  \draw[arr=green] (root) -- (b3);
  \draw[arr=purple] (root) -- (b4);
\end{tikzpicture}
\end{tcolorbox}

\vspace{-1.5mm}
\noindent
\begin{minipage}[t]{0.49\textwidth}\vspace{0pt}
\begin{tcolorbox}[
  enhanced,
  colback=green!4!white,
  colframe=green!60!black,
  arc=2mm, boxrule=0.8pt,
  left=2mm, right=2mm, top=1.8mm, bottom=1.8mm,
  title={\footnotesize\bfseries\color{white}\faCheckCircle\ 4. Ví dụ thực hành chuẩn mực}
]
{\footnotesize
\textbf{1. Tính nhẩm dãy tính:}\\
$8 + 6 - 5 = 14 - 5 = \mathbf{9}$.\\
$17 - 9 + 4 = 8 + 4 = \mathbf{12}$.
\textbf{2. So sánh:}\\
Điền dấu vào chỗ chấm: $9 + 4 \dots 8 + 6$.\\
Ta có: $9 + 4 = 13$; $8 + 6 = 14$.\\
Vì $13 < 14$ nên: $9 + 4 < 8 + 6$.
}
\end{tcolorbox}
\end{minipage}
\hfill
\begin{minipage}[t]{0.49\textwidth}\vspace{0pt}
\begin{tcolorbox}[
  enhanced,
  colback=red!4!white,
  colframe=red!75!black,
  arc=2mm, boxrule=0.8pt,
  left=2mm, right=2mm, top=1.8mm, bottom=1.8mm,
  title={\footnotesize\bfseries\color{white}\faExclamationTriangle\ 5. Mẹo tránh bẫy trắc nghiệm}
]
{\footnotesize
\begin{itemize}\setlength{\itemsep}{1pt}\setlength{\leftskip}{-2.5mm}
  \item \textbf{Bẫy tính vội dãy tính:}\\
  Phải thực hiện từ trái sang phải, không được nhẩm phép trừ trước nếu nó đứng sau!
  \item \textbf{Mẹo cộng ghép 10 khi có 3 số:}\\
  $7 + 8 + 3 = (7 + 3) + 8 = 10 + 8 = 18$.
  \item Khi giải toán có lời văn, chú ý đọc kĩ đối tượng được hỏi là ai.
\end{itemize}
}
\end{tcolorbox}
\end{minipage}


\end{document}
